\begin{thebibliography}{99}
\bibitem{BasuPollackRoy}
S. Basu, R. Pollack and M.-F. Roy, \emph{Algorithms in Real Algebraic Geometry}, second edition, Algorithms and Computation in Mathematics 10, Springer, Berlin, 2006. doi:10.1007/3-540-33099-2.
\bibitem{BewleyKohlberg}
T. Bewley and E. Kohlberg, The asymptotic theory of stochastic games, \emph{Math. Oper. Res.} \textbf{1} (1976), 197--208. doi:10.1287/moor.1.3.197.
\bibitem{BewleyKohlbergStationary}
T. Bewley and E. Kohlberg, On stochastic games with stationary optimal strategies, \emph{Math. Oper. Res.} \textbf{3} (1978), 104--125. doi:10.1287/moor.3.2.104.
\bibitem{Blackwell}
D. Blackwell, Equivalent comparisons of experiments, \emph{Ann. Math. Statist.} \textbf{24} (1953), 265--272. doi:10.1214/aoms/1177729032.
\bibitem{ClarkeCheflesBarnettRiis}
R. B. M. Clarke, A. Chefles, S. M. Barnett and E. Riis, Experimental demonstration of optimal unambiguous state discrimination, \emph{Phys. Rev. A} \textbf{63} (2001), 040305. doi:10.1103/PhysRevA.63.040305.
\bibitem{DenardoFox}
E. V. Denardo and B. L. Fox, Multichain Markov renewal programs, \emph{SIAM J. Appl. Math.} \textbf{16} (1968), 468--487. doi:10.1137/0116038.
\bibitem{FedergrunSchweitzer}
A. Federgrun and P. J. Schweitzer, A fixed point approach to undiscounted Markov renewal programs, \emph{SIAM J. Algebraic Discrete Methods} \textbf{5} (1984), 539--550. doi:10.1137/0605052.
\bibitem{Fritz}
T. Fritz, A synthetic approach to Markov kernels, conditional independence and theorems on sufficient statistics, \emph{Adv. Math.} \textbf{370} (2020), 107239. doi:10.1016/j.aim.2020.107239.
\bibitem{GrafLuschgy}
S. Graf and H. Luschgy, \emph{Foundations of Quantization for Probability Distributions}, Lecture Notes in Mathematics 1730, Springer, Berlin, 2000. doi:10.1007/BFb0103945.
\bibitem{HellmanCover}
M. E. Hellman and T. M. Cover, Learning with finite memory, \emph{Ann. Math. Statist.} \textbf{41} (1970), 765--782. doi:10.1214/aoms/1177696958.
\bibitem{LamondPuterman}
B. F. Lamond and M. L. Puterman, Generalized inverses in discrete time Markov decision processes, \emph{SIAM J. Matrix Anal. Appl.} \textbf{10} (1989), 118--134. doi:10.1137/0610009.
\bibitem{Meyer}
C. D. Meyer, Jr., The role of the group generalized inverse in the theory of finite Markov chains, \emph{SIAM Rev.} \textbf{17} (1975), 443--464. doi:10.1137/1017044.
\bibitem{ShaliziCrutchfield}
C. R. Shalizi and J. P. Crutchfield, Computational mechanics: Pattern and prediction, structure and simplicity, \emph{J. Stat. Phys.} \textbf{104} (2001), 817--879. doi:10.1023/A:1010388907793.
\bibitem{Shaw}
S.-Y. Shaw, Convergence of ergodic nets and approximate solutions of linear functional equations, \emph{RIMS K\^oky\^uroku} \textbf{1415} (2005), 29--37.
\bibitem{GlynnInfanger}
P. W. Glynn and A. Infanger, Solution representations for Poisson's equation, martingale structure, and the Markov chain central limit theorem, \emph{Stochastic Systems} \textbf{14} (2024), 47--68; published online 28 December 2023. doi:10.1287/stsy.2022.0001.
\end{thebibliography}
